An evaluation is a lens through which we view the capabilities of a model. As
language models become increasingly capable at programming, mathematics, and
logical reasoning, the design of this lens becomes central to scientific
progress: a benchmark can reveal a capability that was previously hidden, but it
can also distort our understanding when it measures the wrong behavior or
contains confounding biases.

This thesis studies three lessons from evaluating language models on reasoning
tasks. First, evaluations must focus on the right metrics without bias. For code
and mathematical reasoning, final correctness alone can miss important
distinctions: a model may pass tests without understanding code, solve an
inequality while producing an invalid proof, or perform well on a contaminated
benchmark without genuinely generalizing. Second, benchmark scores are only a
starting point for understanding model behavior. Deeper analysis can reveal
structure hidden by aggregate performance, such as differences between code
generation and code execution, benchmark-specific overfitting, and complementary
failure modes between chain-of-thought prompting and neurosymbolic reasoning.
Third, what we measure shapes what we build. Once evaluations become standard,
they define what counts as progress and influence the capabilities that future
systems are optimized to demonstrate.

Together, these lessons argue that evaluation should be treated not as a
secondary measurement step, but as a central part of building and understanding
AI systems for code and reasoning.
